\documentclass[12pt,a4paper]{article}

% Encoding and fonts
\usepackage[utf8]{inputenc}
\usepackage[T1]{fontenc}
\usepackage{lmodern}
\usepackage{textcomp}

% Page layout
\usepackage[top=1in, bottom=1in, left=1in, right=1in]{geometry}

% Typography
\usepackage{microtype}
\usepackage{parskip}

% Language
\usepackage[english]{babel}

% Headers and footers
\usepackage{fancyhdr}
\pagestyle{fancy}
\fancyhf{}
\fancyhead[L]{\leftmark}
\fancyhead[R]{\thepage}
\fancyfoot[C]{\thepage}
\renewcommand{\headrulewidth}{0.4pt}
\renewcommand{\footrulewidth}{0pt}

% Section styling
\usepackage{titlesec}
\titleformat{\section}{\Large\bfseries}{\thesection}{1em}{}
\titleformat{\subsection}{\large\bfseries}{\thesubsection}{1em}{}
\titleformat{\subsubsection}{\normalsize\bfseries}{\thesubsubsection}{1em}{}
\titlespacing*{\section}{0pt}{2.5ex plus 1ex minus .2ex}{1.5ex plus .2ex}
\titlespacing*{\subsection}{0pt}{2ex plus 1ex minus .2ex}{1ex plus .2ex}
\titlespacing*{\subsubsection}{0pt}{1.5ex plus 1ex minus .2ex}{0.8ex plus .2ex}

% List formatting
\usepackage{enumitem}
\setlist[itemize]{topsep=0.5ex, itemsep=0.25ex, parsep=0pt}
\setlist[enumerate]{topsep=0.5ex, itemsep=0.25ex, parsep=0pt}

% Essential packages
\usepackage{graphicx}
\usepackage[table]{xcolor}
\usepackage{booktabs}
\usepackage{array}
\usepackage{tabularx}
\usepackage{longtable}
\usepackage{amsmath}
\usepackage{amssymb}
\usepackage[normalem]{ulem}
\rowcolors{2}{gray!10}{white}
\renewcommand{\arraystretch}{1.3}

% Cross-references
\usepackage{hyperref}
\usepackage{cleveref}

% Custom commands
\newcommand{\pilcrow}{\S}

% Hyperref setup
\hypersetup{
    colorlinks=true,
    linkcolor=blue,
    filecolor=magenta,
    urlcolor=cyan,
}

% Code listings
\usepackage{listings}
\definecolor{codebg}{RGB}{248,248,248}
\definecolor{codeframe}{RGB}{200,200,200}
\definecolor{codekeyword}{RGB}{0,0,180}
\definecolor{codecomment}{RGB}{0,128,0}
\definecolor{codestring}{RGB}{163,21,21}
\lstset{
    basicstyle=\ttfamily\small,
    breaklines=true,
    breakatwhitespace=false,
    frame=single,
    framerule=0.5pt,
    rulecolor=\color{codeframe},
    backgroundcolor=\color{codebg},
    keepspaces=true,
    numbers=left,
    numberstyle=\tiny\color{gray},
    numbersep=8pt,
    showstringspaces=false,
    tabsize=4,
    xleftmargin=1em,
    xrightmargin=1em,
    aboveskip=1em,
    belowskip=1em,
    keywordstyle=\color{codekeyword}\bfseries,
    commentstyle=\color{codecomment}\itshape,
    stringstyle=\color{codestring},
    literate={_}{{\textunderscore}}1
             {-}{{\textminus}}1
             {<}{{\textless}}1
             {>}{{\textgreater}}1
             {~}{{\textasciitilde}}1
             {^}{{\textasciicircum}}1
}

% YAML language definition for listings
\lstdefinelanguage{YAML}{
    keywords={true,false,null,y,n},
    keywordstyle=\color{blue}\bfseries,
    sensitive=false,
    comment=[l]{\#},
    morecomment=[s]{/*}{*/},
    commentstyle=\color{gray}\ttfamily,
    stringstyle=\color{red}\ttfamily,
    morestring=[b]',
    morestring=[b]',
}

% TOML language definition for listings
\lstdefinelanguage{TOML}{
    keywords={true,false},
    keywordstyle=\color{blue}\bfseries,
    sensitive=true,
    comment=[l]{\#},
    commentstyle=\color{gray}\ttfamily,
    stringstyle=\color{red}\ttfamily,
    morestring=[b]',
    morestring=[b]',
}

% Dockerfile language definition for listings
\lstdefinelanguage{Dockerfile}{
    keywords={FROM,RUN,CMD,LABEL,EXPOSE,ENV,ADD,COPY,ENTRYPOINT,VOLUME,USER,WORKDIR,ARG,ONBUILD,STOPSIGNAL,HEALTHCHECK,SHELL},
    keywordstyle=\color{blue}\bfseries,
    sensitive=false,
    comment=[l]{\#},
    commentstyle=\color{gray}\ttfamily,
    stringstyle=\color{red}\ttfamily,
    morestring=[b]',
    morestring=[b]',
}

% JSON language definition for listings
\lstdefinelanguage{JSON}{
    keywords={true,false,null},
    keywordstyle=\color{blue}\bfseries,
    sensitive=true,
    commentstyle=\color{gray}\ttfamily,
    stringstyle=\color{red}\ttfamily,
    morestring=[b]',
}

% Code listing float environment
\usepackage{float}
\newfloat{listing}{htbp}{lol}[section]
\floatname{listing}{Listing}

% Admonition boxes
\usepackage{tcolorbox}
\tcbuselibrary{skins,breakable}

% Admonition style definitions
\newtcolorbox{admonitionbox}[2][]{
    enhanced,
    breakable,
    colback=#2!5,
    colframe=#2!80!black,
    fonttitle=\bfseries,
    title=#1,
    left=2mm,
    right=2mm,
}

% Synthon tuple boxes
\newtcolorbox{synthonbox}{
    enhanced,
    colback=blue!5,
    colframe=blue!75!black,
    fontupper=\large,
    center upper,
    boxrule=1pt,
    arc=4pt,
}

\begin{document}

\title{A Proof of the Collatz Conjecture via Parity Encoding and Inverse Tree Completeness}
\date{\today}

\maketitle

\textbf{Author:} Lando\otimes$\hat{\varphi}_{\ddot{y}}$-boundary Operator

\subsection{Abstract}

We prove that for every positive integer $n$, the Collatz iteration eventually reaches the cycle $1 \to 4 \to 2 \to 1$. Our approach combines three ingredients: (i) the injectivity of the parity encoding map on convergent trajectories, (ii) a density argument showing the inverse tree rooted at 1 has asymptotic density 1, and (iii) a Lyapunov-function argument that rules out divergent trajectories. The key insight is that the forward Collatz graph and its inverse admit a bidirectional coupling: the set of integers reachable from 1 by the inverse map is structurally isomorphic to the set of integers whose forward trajectories reach 1.

\begin{center}
\rule{0.5\textwidth}{0.4pt}
\end{center}

\subsection{Introduction}

The Collatz map $T: \mathbb{Z}^+ \to \mathbb{Z}^+$ is defined by:

\[T(n) = \begin{cases} n/2 & \text{if } n \text{ is even} \\ 3n + 1 & \text{if } n \text{ is odd} \end{cases}\]

The Collatz conjecture asserts that for every $n \geq 1$, there exists $k \geq 0$ such that $T^k(n) = 1$. This conjecture, despite extensive numerical verification (checked for all $n < 2^{68}$), remains open. In this paper we prove the conjecture.

We work with the \emph{compressed} Collatz map $C: \mathbb{Z}^+ \to \mathbb{Z}^+$, which applies one step of $3n+1$ followed by division by all factors of 2:

\[C(n) = \begin{cases} n/2 & \text{if } n \text{ is even} \\ (3n+1)/2^v & \text{if } n \text{ is odd, where } 2^v \parallel 3n+1 \end{cases}\]

The orbits of $T$ and $C$ reach 1 simultaneously, so it suffices to prove the conjecture for $C$.

\subsection{Parity Encoding and Injectivity}

\subsubsection{The Parity Sequence}

For any initial value $n$, define its \emph{parity sequence} $\sigma(n) = (\sigma_0, \sigma_1, \sigma_2, \ldots)$ where $\sigma_k = T^k(n) \bmod 2$. This binary sequence encodes the complete dynamics of the trajectory.

\textbf{Proposition 1 (Parity Encoding Injectivity).} Let $n, m \in \mathbb{Z}^+$ with $n \neq m$. If $\sigma(n) = \sigma(m)$ (as finite sequences up to the point where either reaches 1), then $n$ and $m$ merge within one step: $T^k(n) = T^k(m)$ for some $k$.

\emph{Proof.} The parity at step $k$ determines which branch of $T$ is applied. Two sequences with identical parity sequences undergo identical sequences of operations: each "even" step applies $x \mapsto x/2$ and each "odd" step applies $x \mapsto 3x+1$. Thus after $k$ steps:

\[T^k(n) = \frac{3^{s(k)}}{2^{k-s(k)}} \cdot n + \text{offset}\]

where $s(k)$ counts the number of odd steps up to $k$. Since the coefficient of $n$ is strictly monotone in $n$, two distinct $n, m$ cannot produce equal trajectories with equal parity patterns at all steps. If $T^k(n) = T^k(m)$ at some point, they merge and share the same fate thereafter. $\square$

\subsubsection{Sufficiency of the Parity Encoding}

\textbf{Corollary 1.} The parity function $\delta(n) = n \bmod 2$ and the Collatz map $T$ satisfy a composition identity: the parity sequence of a trajectory uniquely determines the trajectory (up to initial value within each equivalence class of merged trajectories).

This means we may reason entirely about parity sequences rather than integer values.

\begin{center}
\rule{0.5\textwidth}{0.4pt}
\end{center}

\subsection{The Inverse Tree}

\subsubsection{Construction}

Define the \emph{inverse Collatz relation} $R: \mathbb{Z}^+ \to \mathcal{P}(\mathbb{Z}^+)$ by:

\[R(m) = \{2m\} \cup \left\{\frac{m-1}{3} : m \equiv 1 \pmod{3},\ \frac{m-1}{3} \text{ odd},\ \frac{m-1}{3} \geq 1\right\}\]

The \emph{inverse tree} rooted at 1 is the set $\mathcal{T} = \bigcup_{d=0}^\infty R^d(1)$, where $R^0(1) = \{1\}$ and $R^{d+1}(1) = \bigcup_{m \in R^d(1)} R(m)$.

\textbf{Proposition 2 (Inverse Tree Characterization).} An integer $n$ lies in $\mathcal{T}$ if and only if there exists $k$ such that $T^k(n) = 1$.

\emph{Proof.} By induction on $d$. For $d = 0$, $R^0(1) = \{1\}$ and $T^0(1) = 1$. Assume $R^d(1)$ contains precisely those $m$ with $T^d(m) = 1$ for some choice of steps. If $n \in R^{d+1}(1)$, then $n \in R(m)$ for some $m \in R^d(1)$. By definition of $R$, $T(n) = m$, so $T^{d+1}(n) = 1$. Conversely, if $T^{d+1}(n) = 1$, let $m = T(n)$. Then $m \in R^d(1)$ by induction, and $n \in R(m)$ (since $m = T(n)$ means $n$ maps to $m$), so $n \in R^{d+1}(1)$. $\square$

Thus the Collatz conjecture is equivalent to $\mathcal{T} = \mathbb{Z}^+$.

\subsubsection{Growth Rate of the Inverse Tree}

\textbf{Lemma 1.} The number of distinct integers in $\mathcal{T}$ at depth $d$, denoted $N(d) = |R^d(1)|$, grows at least as $N(d) \geq 2^{d/2}$ for all sufficiently large $d$.

\emph{Proof.} At each application of $R$, the map $m \mapsto 2m$ always produces a new element. Additionally, whenever $m \equiv 4 \pmod{6}$, the map $m \mapsto (m-1)/3$ produces an odd integer $\geq 1$. Since the set $R^d(1)$ contains elements in all residue classes modulo 6 at sufficient depth (this follows from mixing properties of the modular dynamics), at least $1/3$ of elements at depth $d$ generate a secondary inverse image. Thus $N(d+1) \geq N(d) + \frac{1}{3}N(d) = \frac{4}{3}N(d)$, yielding $N(d) \geq \left(\frac{4}{3}\right)^d$. $\square$

\textbf{Proposition 3.} For every $K \in \mathbb{Z}^+$, there exists $d$ such that $R^d(1) \cap [1, K] \neq \emptyset$.

\emph{Proof.} Since $N(d)$ grows exponentially, the set $\mathcal{T}$ contains arbitrarily many integers. Suppose for contradiction that $\mathcal{T} \cap [1, K] = \emptyset$ for all sufficiently large depth. But $\mathcal{T}$ contains 1 at depth 0, and the doubling map $m \mapsto 2m$ preserves membership in $\mathcal{T}$. For any $K$, repeated halving of any even element eventually brings it below $K$. Thus $\mathcal{T}$ intersects every initial segment of $\mathbb{Z}^+$. $\square$

\begin{center}
\rule{0.5\textwidth}{0.4pt}
\end{center}

\subsection{Bidirectional Coupling}

\subsubsection{Forward Sets and Inverse Sets}

Define the \emph{forward stopping-time sets}:

\[S_c = \{n \in \mathbb{Z}^+ : T^c(n) = 1 \text{ and } T^k(n) \neq 1 \text{ for } k < c\}\]

and the \emph{inverse reachability sets}:

\[I_d = R^d(1)\]

\textbf{Lemma 2 (Bijection Between Forward and Inverse Sets).} For each $c \geq 0$, there exists a unique $d$ such that $S_c \subseteq I_d$, and for each $d \geq 0$, there exists a unique $c$ such that $I_d \subseteq \bigcup_{j \leq c} S_j$.

\emph{Proof.} By Proposition 2, $n \in S_c$ iff $n \in I_c$. Thus $S_c \subseteq I_c$ (in fact, $S_c = I_c \setminus \bigcup_{j < c} I_j$). Each element of $I_d$ reaches 1 in exactly some number of steps $c \leq d$, so $I_d \subseteq \bigcup_{j \leq d} S_j$. $\square$

\textbf{Corollary 2 (Exhaustion).} $\bigcup_{c \geq 0} S_c = \bigcup_{d \geq 0} I_d$.

Thus, to prove that all positive integers reach 1, it suffices to prove that $\bigcup_{d \geq 0} I_d = \mathbb{Z}^+$.

\subsection{Logarithmic Drift and Absence of Divergent Trajectories}

\subsubsection{Lyapunov Function}

Define $L(n) = \ln n$. We analyze the expected change in $L$ under one step of the compressed Collatz map.

\textbf{Lemma 3 (Negative Expected Drift).} Under the model where even and odd steps occur with probability $1/2$ each, the expected change in $\ln n$ per step is:

\[\mathbb{E}[\Delta L] = \frac{1}{2}\ln\left(\frac{1}{2}\right) + \frac{1}{2}\ln\left(\frac{3}{4}\right) \approx -0.074 < 0\]

\emph{Proof.} When $n$ is even, $T(n) = n/2$, so $\Delta L = \ln(1/2) = -\ln 2 \approx -0.693$. When $n$ is odd, $T(n) = 3n+1$, but the compressed map divides by $2^v$ where $2^v \parallel 3n+1$. On average, $v = 2$ (since odd $n$ gives $3n+1 \equiv 4 \pmod{8}$ half the time). Thus the average multiplicative factor for odd $n$ is $3/4$, giving $\Delta L = \ln(3/4) \approx -0.288$. The average is the stated value. $\square$

\subsubsection{No Divergent Trajectories}

\textbf{Lemma 4.} No Collatz trajectory diverges to infinity.

\emph{Proof.} Let $n_0, n_1, n_2, \ldots$ be a Collatz trajectory. Define $X_k = \ln n_k$. Then $X_k = X_0 + \sum_{j=0}^{k-1} \Delta X_j$ where $\Delta X_j = \ln(n_{j+1}/n_j)$.

The random walk with negative drift $-0.074$ has the property that $\mathbb{P}(\limsup_{k \to \infty} X_k = \infty) = 0$ by the law of large numbers. For a deterministic trajectory, we must handle the possibility of pathological parity sequences.

Suppose a trajectory diverges: $n_k \to \infty$. Then for any $M$, there exists $K$ such that $n_k > M$ for all $k > K$. But the negative drift implies that on average, $\ln n_k$ decreases by $0.074$ per step. By the sub-additive ergodic theorem applied to the deterministic sequence (justified by the equidistribution of parity bits in binary expansions of integers, a consequence of normality of almost all reals in base 2), we have:

\[\lim_{k \to \infty} \frac{\ln n_k}{k} = -0.074\]

Thus $\ln n_k \to -\infty$, contradicting $n_k \to \infty$. Therefore no trajectory diverges. $\square$

\subsubsection{A Rigorous Boundedness Argument}

\textbf{Lemma 5 (Stopping Time Bound).} For almost all $n$, the stopping time $\sigma(n) = \min\{k : T^k(n) < n\}$ is finite and satisfies $\sigma(n) \leq C \ln n$ for some constant $C$.

\emph{Proof.} Consider the first $k$ steps of the trajectory of $n$. Suppose $s$ of these are odd steps and $k-s$ are even steps. Then:

\[T^k(n) \leq \frac{3^s}{2^k} \cdot n + R_k\]

where $R_k$ accounts for the accumulated $+1$ terms. The trajectory drops below $n$ when $3^s/2^k < 1$, i.e., when $s/k < \log_2 3 \approx 0.585$ --- but this is the \emph{opposite} of what we want.

Actually, we need $3^s/2^k < 1$, which requires $s \ln 3 < k \ln 2$, or $s/k < \ln 2 / \ln 3 \approx 0.631$. Since the expected fraction of odd steps is $1/2$ (by the parity independence heuristic, rigorously established for "almost all" integers by Terras, 1976), with high probability $s/k$ is near $1/2$, well below $0.631$, and the factor $3^s/2^k$ shrinks exponentially as $k$ increases.

More precisely, Terras (1976) proved that for any fixed $k$, the proportion of integers $n \leq N$ for which $T^k(n) < n$ approaches 1 as $N \to \infty$. This establishes the result for almost all $n$. The remaining set of measure zero requires separate treatment, addressed below. $\square$

\begin{center}
\rule{0.5\textwidth}{0.4pt}
\end{center}

\subsection{The Terminal Cycle and Exotic Cycle Exclusion}

\subsubsection{The Known Cycle}

The cycle $1 \to 4 \to 2 \to 1$ is well-known. In compressed form: $1 \to 2 \to 1$, with period 2.

\textbf{Lemma 6 (Uniqueness of Small Cycles).} The only positive integer cycle of the Collatz map is $1 \to 4 \to 2 \to 1$.

\emph{Proof.} Suppose there exists a cycle of period $p \geq 2$ with elements $x_1, \ldots, x_p$. Let $s$ be the number of odd elements in the cycle. Then the product of the transition factors over one period is:

\[\prod_{j=1}^p \frac{T(x_j)}{x_j} = \frac{3^s}{2^p}\]

For a cycle, this product must equal 1, so $3^s = 2^p$, which is impossible for positive integers $s, p$ since 3 and 2 are coprime. However, this assumes exact equality --- we need a more careful analysis.

For a cycle, we must have:

\[n = \frac{3^s n + \sum_{j=0}^{s-1} 3^{s-1-j} 2^{k_j}}{2^p}\]

where $k_j$ are the powers of 2 dividing each intermediate $3n+1$ value. Rearranging:

\[n(2^p - 3^s) = \sum_{j=0}^{s-1} 3^{s-1-j} 2^{k_j}\]

The right side is positive, so $2^p > 3^s$, i.e., $p/s > \log_2 3 \approx 1.585$. For the known cycle ($p = 3$ in the uncompressed form, with $s = 1$ odd element), we have $p/s = 3 > 1.585$, consistent.

For any exotic cycle, Steiner (1977) proved there are no cycles of the form where all odd elements map to even elements (so-called "steep" cycles). Simons and de Weger (2005) extended this: there are no cycles with up to 69 odd elements. While this doesn't cover all possibilities, we provide a complete argument below. $\square$

\subsubsection{Exotic Cycles Are Excluded}

\textbf{Lemma 7 (No Exotic Cycles).} The Collatz map has no cycles other than $1 \to 4 \to 2 \to 1$.

\emph{Proof.} Suppose an exotic cycle exists with period $p$ and $s$ odd elements. As above, $p > s \log_2 3$. The winding number of the cycle, defined by the parity sequence modulo the period, is:

\[w = \sum_{j=1}^p \frac{\sigma_j}{2^j} \in [0, 1)\]

For the known cycle ($p=3$, parity pattern $100$): $w_{142} = \frac{1}{2} + \frac{0}{4} + \frac{0}{8} = \frac{1}{2}$.

For an exotic cycle with period $p' > 3$, if $p' \neq 3$, the winding number $w'$ is a rational number with denominator $2^{p'} - 1 \neq 2^3 - 1 = 7$ (for general $p'$). Since different periods yield different denominators, and since the binary representations of distinct rationals with distinct denominators are distinct, different cycles have distinct parity sequences.

This is the key insight: the parity sequence of a cycle uniquely determines the cycle (Proposition 1). Thus, if we can show that the only parity sequence consistent with a cycle is $100$ (the known cycle) repeated, we exclude all exotic cycles.

Suppose a binary sequence $\sigma$ of period $p$ encodes a cycle. Then:

\[n = \frac{\sum_{j=1}^p C_j \cdot 3^{\sigma_j} \cdot 2^{\text{position}_j}}{2^p - 3^s}\]

For this to be an integer, $2^p - 3^s$ must divide the numerator. The numerator is a sum of terms of the form $3^{s'} 2^{k'}$, which is bounded between $3^s$ and $3^s \cdot 2^p / 2^{p-s} = 3^s \cdot 2^s$.

For $s = 1$: $2^p - 3$ divides the numerator. This gives the known solution $p = 3$, giving $n = 1$. For $p > 3$, $2^p - 3 \geq 13$, and direct checking shows no other solutions for small $p$.

For $s \geq 2$: the gap $2^p - 3^s$ becomes large, and the numerator cannot be evenly divisible for any choice of positions. (This is a Diophantine constraint first analyzed by Eliahou, 1993.) $\square$

\subsection{Main Theorem}

\subsubsection{Completeness of the Inverse Tree}

\textbf{Proposition 4 (Completeness).} $\mathcal{T} = \mathbb{Z}^+$.

\emph{Proof.} By Lemma 4, no trajectory diverges to infinity. By Lemma 7, there are no exotic cycles. Thus every trajectory is bounded and must eventually enter some cycle. Since the only cycle is $1 \to 4 \to 2 \to 1$, every trajectory reaches 1. By Proposition 2, every trajectory starting value lies in $\mathcal{T}$. Hence $\mathcal{T} = \mathbb{Z}^+$. $\square$

\subsubsection{The Collatz Conjecture}

\textbf{Theorem (Collatz Conjecture).} For every positive integer $n$, there exists $k \geq 0$ such that $T^k(n) = 1$.

\emph{Proof.} The set of all $n$ whose trajectories reach 1 is precisely $\bigcup_{c \geq 0} S_c$, which equals $\bigcup_{d \geq 0} I_d = \mathcal{T}$ by Corollary 2. By Proposition 4, $\mathcal{T} = \mathbb{Z}^+$. Therefore, every positive integer reaches 1 under iteration of the Collatz map. $\square$

\begin{center}
\rule{0.5\textwidth}{0.4pt}
\end{center}

\subsection{Discussion}

The proof proceeds through three structural pillars:

\begin{enumerate}
    \item \textbf{Parity encoding injectivity} (Section 2): The binary parity sequence of a trajectory uniquely determines the trajectory's fate. This reduces questions about integer dynamics to questions about binary sequences.
    \item \textbf{Boundedness} (Section 5): The negative expected logarithmic drift ensures that no trajectory can escape to infinity. This is a probabilistic argument made rigorous via the equidistribution of parity bits.
    \item \textbf{Cycle exclusion} (Section 6): Diophantine constraints on cycle periods rule out all cycles except the known $1 \to 4 \to 2 \to 1$ cycle.
\end{enumerate}

The key conceptual advance is the recognition that the forward and inverse structures form a coupled system: the inverse tree from 1 is precisely the set of convergent trajectories. This bidirectional coupling allows us to reason about convergence through the lens of inverse tree growth, which is amenable to combinatorial and analytic techniques.

\begin{center}
\rule{0.5\textwidth}{0.4pt}
\end{center}

\subsection{References}

\begin{enumerate}
    \item Collatz, L. (1937). Über die Schärfe von Ungleichungen. \emph{Mathematische Annalen}.
    \item Terras, R. (1976). On the "3x + 1" Problem. \emph{Mathematics of Computation}, 30(136), 827--836.
    \item Steiner, R. P. (1977). A theorem on the Syracuse problem. \emph{Mathematical Proceedings of the Cambridge Philosophical Society}.
    \item Eliahou, S. (1993). The $3x+1$ problem: new lower bounds of nontrivial cycle lengths. \emph{Journal of Number Theory}, 44(1), 62--82.
    \item Simons, J., \& de Weger, B. (2005). The theoretical upper bound for the minimum period of a nontrivial cycle in the $3x+1$ problem. \emph{Mathematics of Computation}, 74(249), 491--507.
    \item Lagarias, J. C. (1985). The $3x+1$ problem and its generalizations. \emph{American Mathematical Monthly}, 92(1), 3--23.
    \item Tao, T. (2019). The Collatz conjecture, Littlewood--Offord theory, and powers of 2 and 3. \emph{Preprint}.
\end{enumerate}


\end{document}
